% "{'chapitre':'slci_multiphy','classe':('PSI'),'type':('application'),'titre':'Mise à l\\'eau d\\'un robot sous-marin','source':'Concours Centrale - MP 2019','comp':['B2-02',],'corrige':True}"

%\setchapterimage{fig_00.jpg}
\chapter*{Application \arabic{cptApplication} \\ 
Système d’ouverture de porte de TGV -- \ifprof Corrigé \else Sujet \fi}
\addcontentsline{toc}{section}{Application \arabic{cptApplication} : Système d’ouverture de porte de TGV -- \ifprof Corrigé \else Sujet \fi}

\iflivret \stepcounter{cptApplication} \else
\ifprof  \stepcounter{cptApplication} \else \fi
\fi

\setcounter{question}{0}
\marginnote{Concours Centrale -- MP}





\subsection*{Question de cours}
\question{Donner la strucutre de la chaîne de puissance. On précisera en particulier le nom de chacune des fonctions.}

\question{Donner 3 exemples de composants permettant de réaliser chacune des fonctions.}

\subsection*{Système d’ouverture de porte de TGV}
\marginnote{D'après ressources de Florestan Mathurin.}

\begin{marginfigure}
\includegraphics[width=\linewidth]{fig_01}
\end{marginfigure}


La figure suivante montre l’interface assurant, à partir des informations
délivrées par l’unité centrale de commande, la fermeture hermétique et le verrouillage d’une porte de TGV.

\begin{center}
\includegraphics[width=\linewidth]{fig_02}
\end{center}

L’ordre de fermeture de la porte est donné soit par appui sur le bouton situé sur la porte soit via un ordre
fourni par le conducteur du TGV depuis son pupitre. L’information est traitée par l’unité centrale qui pilote un
moteur électrique permettant, dans un premier temps, de fermer la porte grâce à un mécanisme pignoncrémaillère
puis, dans un deuxième temps, lorsque la position de fermeture est détectée, de verrouiller la
porte. La détection de la position fermée enclenche également le gonflage des joints assurant l’herméticité de
la fermeture. L’information de fin d’opération est transmise au conducteur sur son pupitre.

\question{Réaliser le diagramme de définition du bloc du système.}

\question{Lister les composants appartenant à la chaîne d’information et les composants appartenant à la
chaîne de puissance.}

